Chương này trình bày các cơ sở lý thuyết và công nghệ nền tảng phục vụ quá trình thiết kế hệ thống. Nội dung được triển khai theo hai nhóm chính: cơ sở về chất lượng không khí trong nhà và điều kiện chăm sóc cây lan ý; cơ sở công nghệ gồm cảm biến, truyền thông LoRa, vi điều khiển Arduino Pro Mini, gateway ESP32, nền tảng ThingsBoard và ứng dụng Flutter.

\section{Chất lượng không khí trong nhà}
Chất lượng không khí trong nhà (Indoor Air Quality - IAQ) phản ánh mức độ phù hợp của môi trường không khí đối với sức khỏe, sự thoải mái và hiệu quả sinh hoạt của con người. Trong căn hộ chung cư, IAQ chịu tác động đồng thời của thông gió, số người trong phòng, hoạt động nấu nướng, thiết bị điện, vật liệu nội thất, nhiệt độ, độ ẩm và thói quen sử dụng điều hòa. Do các yếu tố này thay đổi theo thời gian, việc giám sát liên tục có ý nghĩa hơn so với đánh giá cảm quan tại một thời điểm.

Trong phạm vi khóa luận, nhóm thông số môi trường được giám sát gồm nhiệt độ, độ ẩm, nồng độ CO2 và ánh sáng. Nhiệt độ và độ ẩm vừa ảnh hưởng đến cảm giác tiện nghi của người sử dụng, vừa liên quan đến khả năng sinh trưởng của lan ý. CO2 được sử dụng như chỉ báo gián tiếp cho mức độ thông gió và mật độ người trong phòng; khi CO2 tăng cao, hệ thống có thể phát cảnh báo hoặc đề xuất bật quạt thông gió. Ánh sáng được xem xét ở cả hai khía cạnh: điều kiện sinh hoạt của con người và điều kiện quang hợp của cây trồng trong nhà.

\subsection{Nhiệt độ và độ ẩm}
Nhiệt độ quá cao hoặc quá thấp đều có thể gây khó chịu và làm tăng nhu cầu sử dụng điều hòa. Độ ẩm thấp thường gây khô da, khô đường hô hấp, trong khi độ ẩm quá cao tạo điều kiện cho nấm mốc phát triển. Đối với lan ý, tài liệu chăm sóc khuyến nghị môi trường ấm, khoảng 65--80~$^\circ$F, độ ẩm từ khoảng 50\% trở lên và tránh các luồng gió nóng hoặc lạnh thổi trực tiếp vào cây \cite{peace_lily_care_bhg}. Do đó, nhiệt độ và độ ẩm không chỉ là thông số tiện nghi mà còn là cơ sở để đánh giá điều kiện chăm sóc cây trong hệ thống.

\subsection{Nồng độ CO2}
Trong không gian căn hộ, CO2 chủ yếu phát sinh từ hoạt động hô hấp của con người. Khi phòng kín hoặc thông gió kém, nồng độ CO2 có xu hướng tăng và phản ánh mức độ trao đổi không khí với bên ngoài. Các tài liệu về IAQ xem CO2 là một chỉ báo quan trọng đối với ô nhiễm có nguồn gốc từ người sử dụng; từ khoảng 1000 ppm, cảm giác khó chịu, mệt mỏi hoặc giảm tập trung có thể xuất hiện ở một bộ phận người trong phòng \cite{co2_monitoring_rehva}. Vì vậy, hệ thống sử dụng CO2 như một ngưỡng cảnh báo để đề xuất mở cửa, bật quạt thông gió hoặc điều chỉnh cách sử dụng phòng.

\subsection{Hợp chất bay hơi và giới hạn của hệ thống}
Ngoài CO2, IAQ còn chịu ảnh hưởng bởi bụi mịn và các hợp chất hữu cơ bay hơi (VOC) phát sinh từ vật liệu nội thất, sơn, chất tẩy rửa hoặc hoạt động nấu nướng \cite{epa_iaq}. Một số nghiên cứu về cây trồng trong nhà cho thấy cây có thể hỗ trợ hấp thụ một phần chất ô nhiễm trong điều kiện thí nghiệm \cite{nasa_clean_air}. Đối với lan ý, tài liệu về hệ thủy canh cho IAQ đề cập \textit{Spathiphyllum} như loài cây có tiềm năng hỗ trợ hấp thụ CO2 và VOC khi được duy trì trong điều kiện sinh trưởng phù hợp \cite{hydroponic_iaq_survey}. Tuy nhiên, nguyên mẫu của khóa luận chưa đo trực tiếp VOC hoặc bụi mịn; vì vậy, hệ thống tập trung vào CO2, nhiệt độ, độ ẩm, ánh sáng và các thông số chăm sóc cây như nhóm chỉ báo ban đầu cho môi trường căn hộ.
